\section{Tokenization 附录：把课本上的概念落到实现}

\subsection{BPE 的训练和使用}

BPE 的训练过程是一个标准的频次合并循环：从 bytes 开始，统计相邻 pair 的出现次数，选最常见 pair 合并为新 token，重复若干轮，直到达到词表预算。编码时则把字符串逐步压成这些合并结果。

\begin{figure}[H]
\centering
\begin{tikzpicture}[node distance=0.9cm, every node/.style={font=\small, align=center}]
\node[draw, rounded corners, fill=orange!12, minimum width=3.1cm, minimum height=0.8cm] (t1) {start from bytes};
\node[draw, rounded corners, fill=orange!20, below=of t1, minimum width=3.1cm, minimum height=0.8cm] (t2) {count pair frequencies};
\node[draw, rounded corners, fill=orange!28, below=of t2, minimum width=3.1cm, minimum height=0.8cm] (t3) {merge the most common pair};
\node[draw, rounded corners, fill=orange!36, below=of t3, minimum width=3.1cm, minimum height=0.8cm] (t4) {repeat until budget};
\draw[-{Latex[length=3mm]}, thick] (t1) -- (t2);
\draw[-{Latex[length=3mm]}, thick] (t2) -- (t3);
\draw[-{Latex[length=3mm]}, thick] (t3) -- (t4);
\end{tikzpicture}
\caption{BPE 的训练循环：统计、合并、更新、重复}
\end{figure}

\begin{figure}[H]
\centering
\begin{tikzpicture}[every node/.style={font=\small, align=center}]
\node[draw, rounded corners, fill=green!10, minimum width=2.0cm, minimum height=0.8cm] (b1) {b};
\node[draw, rounded corners, fill=green!10, right=0.3cm of b1, minimum width=2.0cm, minimum height=0.8cm] (b2) {o};
\node[draw, rounded corners, fill=green!10, right=0.3cm of b2, minimum width=2.0cm, minimum height=0.8cm] (b3) {o};
\node[draw, rounded corners, fill=green!10, right=0.3cm of b3, minimum width=2.0cm, minimum height=0.8cm] (b4) {k};
\node[draw, rounded corners, fill=yellow!20, below=1.2cm of b2, minimum width=2.8cm, minimum height=0.8cm] (m1) {bo};
\node[draw, rounded corners, fill=yellow!28, right=0.6cm of m1, minimum width=2.8cm, minimum height=0.8cm] (m2) {ok};
\draw[-{Latex[length=3mm]}, thick] (b1) -- (m1);
\draw[-{Latex[length=3mm]}, thick] (b2) -- (m1);
\draw[-{Latex[length=3mm]}, thick] (b3) -- (m2);
\draw[-{Latex[length=3mm]}, thick] (b4) -- (m2);
\end{tikzpicture}
\caption{BPE 的 merge 直觉：高频相邻片段会逐步合并成更大的 token}
\end{figure}

\begin{table}[H]
\centering
\begin{tabular}{p{3cm}p{4cm}p{5.8cm}}
\toprule
\textbf{方案} & \textbf{优点} & \textbf{缺点} \\
\midrule
Character-based & 可逆、直观 & vocab 太大，稀有字符浪费 \\
Byte-based & vocab 很小 & 序列太长，attention 成本高 \\
Word-based & 语义清楚 & OOV 和词表爆炸 \\
BPE & 兼顾压缩率和覆盖率 & 需要训练和高效实现 \\
\bottomrule
\end{tabular}
\caption{几种 tokenization 方案的折中}
\end{table}

\begin{warningbox}{tokenization 的工程边界}
tokenization 不是“越精致越好”，而是越能压缩常见模式、同时不把序列长度拉爆越好。
\end{warningbox}

